\section*{अध्याय \ref{ch:b2:neurons-synapses} --- न्यूरॉन, क्रिया विभव और तंत्रिका-संधियाँ}

\begin{solution}{exo:b2:neurons-synapses:1}
\omterm{def:b2:neurons-synapses:neuron}{द्रुमिकाएँ} \omterm{def:b2:neurons-synapses:synapse}{तंत्रिका-संधियाँ} पाती हैं; सोमा केंद्रक रखता है और
समाकलित करता है; तंत्रिकाक्ष-टीला तय करता है; \omterm{def:b2:neurons-synapses:neuron}{तंत्रिकाक्ष} चालन करता है; और
सिरे संचारी छोड़ते हैं। \omterm{def:b2:neurons-synapses:neuron}{मायलिन} किसी ग्लिया-कोशिका की लपेटी हुई झिल्ली है,
जो कोशिकाद्रव से लगभग मुक्त है; वह \omterm{def:b2:neurons-synapses:neuron}{तंत्रिकाक्ष} को रोधित करता है ताकि
वह आवेग गाँठों के बीच कूदे और कम लागत पर सौ गुना तेज़
चले।
\end{solution}

\begin{solution}{exo:b2:neurons-synapses:2}
देहली तक विध्रुवण; चढ़ती अवस्था --- वोल्टता-द्वारित सोडियम
वाहिकाएँ खुलती हैं, धन प्रतिपुष्टि; $E_{\mathrm{Na}}$ के पास शिखर ---
सोडियम वाहिकाएँ निष्क्रिय होती हैं; गिरती अवस्था --- वोल्टता-द्वारित पोटैशियम
वाहिकाएँ खुलती हैं; अधोक्षेप --- पोटैशियम वाहिकाएँ अब भी खुली,
$E_{\mathrm K}$ के पास; और उनके बंद होते ही उबार। सब-या-कुछ-नहीं क्योंकि वह
धन प्रतिपुष्टि या तो पूरी चलती है या बिलकुल नहीं; और एकतरफ़ा
क्योंकि अभी-अभी पार की गई झिल्ली अपवर्तक है (सोडियम वाहिकाएँ
निष्क्रिय)।
\end{solution}

\begin{solution}{exo:b2:neurons-synapses:3}
आवेग उस सिरे तक पहुँचता है; वोल्टता-द्वारित कैल्सियम वाहिकाएँ खुलती हैं;
कैल्सियम पुटिका-संलयन छेड़ती है; संचारी उस दरार के आरपार
विसरित होता है; पश्चसंधि ग्राहियों से बँधता है; वाहिकाएँ खुलती हैं (या कोई \omterm{prop:b2:cell-signalling:gpcr}{G प्रोटीन}
स्विच करता है); \omterm{def:b2:neurons-synapses:synapse}{पश्चसंधि विभव}; संचारी एंजाइम या
वाहक से हटाया जाता है; पुटिका-झिल्ली वापस ली जाती है।
\end{solution}

\begin{solution}{exo:b2:neurons-synapses:4}
विद्युत्: कोई देरी नहीं, आम तौर पर दोनों दिशाएँ, कोई प्रवर्धन नहीं, कोई
संदमन नहीं (केवल धारा गुज़रती है)। रासायनिक: आधा मिलीसेकंड,
एकतरफ़ा, प्रवर्धन (एक आवेग सैकड़ों अंश छोड़ता है),
संदमन संभव (क्लोराइड या पोटैशियम के लिए वाहिकाएँ), और
परिवर्तनीय।
\end{solution}

\begin{solution}{exo:b2:neurons-synapses:5}
$E = (61.5/z)\log_{10}(c_{\text{बाहर}}/c_{\text{भीतर}})$ से ये मान
मिलते हैं। पोटैशियम \qty{-89}{mV}; सोडियम \qty{+61}{mV}; क्लोराइड, $z =
-1$ के साथ, $-61.5\log_{10} 11 = \qty{-64}{mV}$; कैल्सियम, $z = 2$ के साथ,
$30.75\times 4.3 = \qty{+132}{mV}$। \qty{-70}{mV} पर:
पोटैशियम निकलता है, सोडियम घुसता है, क्लोराइड थोड़ा निकलता है (वह
झिल्ली अपने साम्य से नीचे है), और कैल्सियम ज़ोर से घुसती है।
\end{solution}

\begin{solution}{exo:b2:neurons-synapses:6}
$20 : 1$: $(20\times(-89) + 61)/21 = \qty{-82}{mV}$, यानी $E_{\mathrm K}$ के पास
विश्राम-दशा। $1 : 20$: $(-89 + 20\times 61)/21 =
\qty{+54}{mV}$, यानी वह शिखर, $E_{\mathrm{Na}}$ के पास। $50 : 1$: $(50\times
(-89) + 61)/51 = \qty{-86}{mV}$, यानी अतिरिक्त पोटैशियम वाहिकाएँ खुली रहते
वह अधोक्षेप।
\end{solution}

\begin{solution}{exo:b2:neurons-synapses:7}
$\sqrt{500} = 22$: \qty{22}{m/s}। \qty{100}{m/s} तक पहुँचने के लिए किसी
बिना-मायलिन \omterm{def:b2:neurons-synapses:neuron}{तंत्रिकाक्ष} को $d = 10^{4}$ \unit{\micro m} चाहिए होता, यानी कोई
सेंटीमीटर। एक मीटर \omterm{def:b2:neurons-synapses:neuron}{मायलिन} वाले तंतु में \qty{10}{ms} लेता है,
विद्रूप-तंत्रिकाक्ष में \qty{45}{ms}, और किसी \qty{1}{\micro m} तंतु में
पूरा सेकंड।
\end{solution}

\begin{solution}{exo:b2:neurons-synapses:8}
$m = \ln(200/74) = 1.0$; $P(1) = e^{-1} = 0.37$, $P(2) = 0.18$, $P(3)
= 0.06$। माध्य अनुक्रिया $1.0\times 0.5 = \qty{0.5}{mV}$।
\end{solution}

\begin{solution}{exo:b2:neurons-synapses:9}
$1/\qty{2}{ms} = 500$ आवेग प्रति सेकंड। उस स्पर्श की तीव्रता
आवेगों की बारंबारता में कूटबद्ध है (और इसमें कि कितने \omterm{def:b2:neurons-synapses:neuron}{न्यूरॉन}
भरती किए गए); और वे आवेग स्वयं एक जैसे हैं।
\end{solution}

\begin{solution}{exo:b2:neurons-synapses:10}
$\lambda = \sqrt{R_m d/4R_i}$: \qty{0.5}{mm}, \qty{1.6}{mm},
\qty{11}{mm}। \omterm{def:b2:neurons-synapses:neuron}{मायलिन} के साथ, \qty{10}{\micro m}: $1.6\times\sqrt{200} =
\qty{22}{mm}$। \qty{1}{mm} का कोई अंतर-गाँठ उस संकेत का केवल $1 - e^{-1/22}
= \qty{4}{\%}$ खोता है, इसलिए अगली गाँठ आसानी से देहली तक
ले आई जाती है। किसी बिना-मायलिन \qty{3}{mm} के आरपार वह संकेत गिरकर
$e^{-3/1.6} = \qty{15}{\%}$ रह जाता है --- यानी \qty{100}{mV} के किसी आवेग का, यानी लगभग वह
\qty{15}{mV} देहली: चालन चूक जाता है या भरोसे लायक नहीं रहता।
\end{solution}

\begin{solution}{exo:b2:neurons-synapses:11}
टेट्रोडोटॉक्सिन: कोई सोडियम धारा नहीं, कोई आवेग नहीं, कोई संचरण नहीं।
टेट्राएथिलअमोनियम: कोई पोटैशियम धारा नहीं, लंबा आवेग, प्रति आवेग अधिक
संचारी छूटा। कोई कैल्सियम नहीं: आवेग सामान्य,
कोई मोचन नहीं, कोई संचरण नहीं। क्यूरारे: मोचन सामान्य, ग्राही
रुके, कोई अंत्य-पट्ट विभव नहीं --- यानी लकवा। कोलिनएस्टरेज़
संदमक: ऐसिटिलकोलीन बनी रहती है, ग्राही खुले रहते हैं, वह पेशी
विध्रुवित होकर फिर पुनर्ध्रुवित नहीं हो पाती --- यानी अधिकता से लकवा।
बोटुलिनम विष: पुटिकाएँ जुड़ नहीं सकतीं, कोई मोचन नहीं --- यानी लकवा।
\end{solution}

\begin{solution}{exo:b2:neurons-synapses:12}
वह देरी ख़रीदती है: केवल एक दिशा में संचरण; प्रवर्धन
(कोई एक आवेग सैकड़ों अंश छोड़ता है और किसी बड़ी कोशिका को चला सकता
है); संदमन, जो कोई विद्युत् जोड़ नहीं कर सकता;
और सुनम्यता, क्योंकि अंशों तथा ग्राहियों की संख्या काम में लेने के साथ बदल
सकती है, और परिपथ ऐसे ही सीखते हैं। विद्युत् \omterm{def:b2:neurons-synapses:synapse}{तंत्रिका-संधियाँ} वहाँ काम में ली जाती हैं
जहाँ समकालिकता तथा गति गणना से अधिक मायने रखती हैं:
पलायन-परिपथ, \omterm{def:b2:heart:anatomy}{हृदय} की पेशी, कुछ संदमक जाल।
\end{solution}

\begin{solution}{pb:b2:neurons-synapses:1}
\textbf{1.} $E_{\mathrm K} = 61.5\log_{10}(4/140) = \qty{-95}{mV}$;
$E_{\mathrm{Na}} = 61.5\log_{10}(145/12) = \qty{+67}{mV}$।
\textbf{2.} $(25\times(-95) + 67)/26 = \qty{-89}{mV}$।
\textbf{3.} $E_{\mathrm K} = 61.5\log_{10}(8/140) = \qty{-76}{mV}$;
$V = (25\times(-76) + 67)/26 = \qty{-71}{mV}$: यानी देहली के अधिक पास,
इसलिए पहले-पहल अधिक उत्तेजनीय --- और, यदि वह बना रहे, तो कम, क्योंकि सोडियम
वाहिकाएँ उस विध्रुवित स्तर पर निष्क्रिय हो जाती हैं।
\textbf{4.} वे प्रवणताएँ घंटों भर चुक जाती हैं ज्यों-ज्यों सोडियम भीतर और
पोटैशियम बाहर रिसते हैं; वह विभव शून्य की ओर सरकता है; और पंप के
रुकने पर उस कोशिका के अपारगम्य ऋणायन क्लोराइड तथा जल भीतर खींच लेते हैं,
और वह फूल जाती है।
\textbf{5.} प्रति मिलीमीटर पृष्ठ $\pi\times 15\times 10^{-6}\times
10^{-3} = \qty{4.7e-8}{m^2}$; $C = \qty{4.7e-10}{F}$; $Q = CV =
\qty{4.7e-11}{C}$: यानी $2.9\times 10^{8}$ आयन।
\textbf{6.} प्रति मिलीमीटर आयतन $\pi(7.5\times 10^{-6})^{2}\times
10^{-3} = \qty{1.8e-13}{m^3} = \qty{1.8e-10}{L}$, जो $12\times
10^{-3}\times 1.8\times 10^{-10} = 2.1\times 10^{-12}$ मोल, यानी $1.3\times
10^{12}$ आयन थामे है: यानी $2\times 10^{-4}$ का परिवर्तन।
\textbf{7.} प्रति मिलीमीटर प्रति आवेग $2.9\times 10^{8}/3 \approx 10^{8}$
ATP।
\textbf{8.} हर ओर $0.9/90 = \qty{10}{ms}$।
\textbf{9.} नंगा: $\sqrt{1\times 15\times 10^{-6}/4} = \qty{1.9}{mm}$;
\omterm{def:b2:neurons-synapses:neuron}{मायलिन} वाला: $\times\sqrt{250} = \qty{31}{mm}$।
\textbf{10.} नंगा: $e^{-1.5/1.9} = 0.45$, \qty{45}{mV}; \omterm{def:b2:neurons-synapses:neuron}{मायलिन} वाला:
$e^{-1.5/31} = 0.95$, \qty{95}{mV}। दोनों उस \qty{15}{mV}
देहली से बड़े हैं; पर नंगे \omterm{def:b2:neurons-synapses:neuron}{तंत्रिकाक्ष} को रास्ते भर अपनी पूरी झिल्ली
आवेशित करनी पड़ती है, और वही उसे धीमा बनाता है, जबकि \omterm{def:b2:neurons-synapses:neuron}{मायलिन} के नीचे लगभग
कुछ भी न खोया न आवेशित किया जाता है।
\textbf{11.} $\sqrt{15} = \qty{3.9}{m/s}$: हर ओर \qty{0.23}{s},
यानी पूरे चक्कर के लिए लगभग आधा सेकंड --- यानी ऐसे किसी प्रतिवर्त के लिए बहुत धीमा
जिसे किसी ठोकर को पकड़ना है।
\textbf{12.} $e^{-4/1.9} = 0.12$: दूर की गाँठ पर \qty{12}{mV}, यानी देहली से
नीचे --- वह आवेग रुक जाता है और वह प्रतिवर्त खो जाता है।
\textbf{13.} उस आवेग के पीछे सोडियम वाहिकाएँ मिलीसेकंड-दो मिलीसेकंड तक
निष्क्रिय रहती हैं, इसलिए अभी-अभी पार की गई झिल्ली फिर उत्तेजित नहीं हो
सकती और वह तरंग केवल आगे ही चल सकती है।
\textbf{14.} $m = \ln(300/111) = 1.0$।
\textbf{15.} $P(1) = 0.37$, $P(2) = 0.18$, $P(3) = 0.06$: यानी उन 300
प्रयासों में से लगभग 110, 55 और 18।
\textbf{16.} $1.0\times 0.4 = \qty{0.4}{mV}$।
\textbf{17.} $P(0) = e^{-60} \approx 10^{-26}$: कभी नहीं; माध्य अनुक्रिया
\qty{24}{mV}, यानी उस \qty{15}{mV} देहली से ऊपर: कोई एक ऐसी \omterm{def:b2:neurons-synapses:synapse}{तंत्रिका-संधि}
उस प्रेरक \omterm{def:b2:neurons-synapses:neuron}{न्यूरॉन} को दगा देती है।
\textbf{18.} एक साथ सक्रिय बीस \omterm{def:b2:neurons-synapses:synapse}{तंत्रिका-संधियाँ} अपनी धाराएँ जोड़ती हैं
(स्थानिक संकलन), और अवरैखिक ढंग से ज्यों-ज्यों वह विभव उन वाहिकाओं के
उत्क्रमण विभव के पास आए; हर एक की $m = 1$ पर भी वे
\qty{8}{mV} देतीं, और सामान्य $m$ पर देहली से कहीं अधिक।
\textbf{19.} $E_{\mathrm{Cl}} = V_{\text{विश्राम}}$ पर क्लोराइड वाहिकाएँ खोलना
उस विभव को हिलाए बिना चालकता जोड़ता है; और वह उत्तेजक
धारा अब किसी अधिक रिसती झिल्ली से बहती है और कोई छोटा
विध्रुवण पैदा करती है --- यानी पार्श्वपथ संदमन।
\textbf{20.} $10 + 0.6 + 10 + 0.6 + 5 = \qty{26}{ms}$, मापे गए
\qty{30}{ms} के मुक़ाबले।
\textbf{21.} पेशी-तर्कु में उस ग्राही का अपना सक्रियण, उस तंतु पर
पेशी-तंतु के \omterm{prop:b2:neurons-synapses:ap}{क्रिया विभव} का संचरण
(सेंटीमीटरों पर मीटर प्रति सेकंड), और बल प्रकट होने से पहले
वह यांत्रिक विलंबता।
\textbf{22.} हर \omterm{def:b2:neurons-synapses:synapse}{तंत्रिका-संधि} आधा मिलीसेकंड और कोई अंतरन्यूरॉन एक कोशिका
जोड़ता है; पर अंतरन्यूरॉन उस संकेत को उलटने
(प्रतिपक्षी को संदमित करने), दूसरों से जोड़ने, और
मस्तिष्क से साधे जाने देते हैं --- यानी ऐसा प्रतिवर्त जो गढ़ा जा सके।
\textbf{23.} \qty{1.8}{m} के लिए $90 - 6.2 = \qty{84}{ms}$ चालन:
यानी लगभग \qty{21}{m/s}।
\textbf{24.} हर वोल्टता पर आधी भीतरी धारा: वह देहली
चढ़ती है, वह आवेग छोटा तथा धीमा होता है, हर गाँठ पर वह सुरक्षा-गुणक
गिरता है, और वह आवेग चूक सकता है; वह प्रतिवर्त कमज़ोर होता है या
मिट जाता है।
\textbf{25.} \omterm{thm:b2:neurons-synapses:chord}{विश्राम विभव} \qty{-89}{mV}; हर ओर \qty{10}{ms};
$m = 1.0$; और गिनी गई विलंबता \qty{26}{ms}।
\end{solution}
